\documentclass[letterpaper,11pt]{article}
\usepackage{resume_style}
\begin{document}

%---------- HEADER ----------
\begin{center}
  {\Huge \scshape Vedaang Chopra} \\ \vspace{2pt}
  \small
  \href{mailto:vedaangchopra@gatech.edu}{\underline{vedaangchopra@gatech.edu}} \;|\;
  \underline{+1 (404) 740-9905} \;|\;
  Atlanta, GA \;|\;
  \href{https://linkedin.com/in/vedaang-chopra}{\underline{LinkedIn}} \;|\;
  \href{https://github.com/Vedaang-Chopra}{\underline{GitHub}} \;|\;
  \href{https://vedaangchopra.live/}{\underline{Website}}
  \\ \vspace{3pt}
  \small
  MS CS (ML) at Georgia Tech \;|\; Advised by \textbf{Dr.\ Matthew Gombolay}, CORE Robotics Lab \\
  \textit{Research:} Agentic AI systems with execution-grounded evaluation -- Multimodal model routing \& serving -- Production ML at enterprise scale
\end{center}
\vspace{-6pt}

%---------- RESEARCH EXPERIENCE ----------
\section{Research Experience}
\resumeSubHeadingListStart
\resumeSubheading
{Georgia Institute of Technology -- CORE Robotics Lab @ Siemens}{Atlanta, GA}
{Graduate Research Assistant \;|\; Advisor: \textbf{Dr.\ Matthew Gombolay}}{Jan 2026 -- Present}
\resumeItemListStart
  \resumeItem{Building a closed-loop agentic pipeline for parametric CAD code generation:
  a frozen LLM generates CadQuery programs and execution feedback (compile validity, geometric
  correctness via Chamfer/Hausdorff on STL/B-Rep) drives a LangGraph-orchestrated repair loop --
  inference-time search with verifiable geometric reward; 16-module system maintained as a 202-test automated evaluation harness with VLM-based visual judgment.}

  \resumeItem{Designed scalable data curation and processing for the pipeline's retrieval-augmented
  prompting: structured evidence collection, prompt-evaluation components, and failure classification
  feeding the automated quality-control gates of a 202-test evaluation harness.}
\resumeItemListEnd

\resumeSubheading
{Georgia Institute of Technology -- CS 8903 Special Problems}{Atlanta, GA}
{Graduate Researcher}{Jan 2025 -- Present}
\resumeItemListStart
  \resumeItem{\textbf{ARTEMIS} (Aug 2025--Present): Developing cost-aware routing for Vision-Language
  Models (VLMs): trained neural multi-task router selects the optimal VLM per request across VQA,
  OCR, captioning, and reasoning -- SLA-aware load balancing serves 6+ heterogeneous backends
  (Gemma 3, Qwen-VL, DeepSeek OCR, Llama-4) via vLLM with 5 dynamic routing modes.}

  \resumeItem{\textbf{CERBERUS} (Aug 2025--Present): Encoder-frozen cross-modal alignment research:
  trained alignment stack with PyTorch DDP + mixed precision on H100/A100 clusters; Matryoshka
  Representation Learning compresses embeddings 4096$\to$128 dims at ~96\% retention -- R@5 78\%
  on PixMo vision-text retrieval; ablated Perceiver Resampler integration and diagnosed retrieval
  collapse to initialization/gradient issues.}

  \resumeItem{\textbf{ATHENA} (Jan--May 2025): Built a 5-agent supervisor-worker multi-agent system
  via LangGraph Command routing with reflection-based critique; designed cross-modal evaluation
  framework (BLEU-4, METEOR, BERTScore, CLIPScore) over 100 reference videos -- +0.18 BLEU-4
  absolute over single-agent baseline; ablation cut fact coverage 27\% without the critique agent.}
\resumeItemListEnd
\resumeSubHeadingListEnd

%---------- INDUSTRY EXPERIENCE ----------
\section{Industry Experience}
\resumeSubHeadingListStart
\resumeSubheading
{Fortinet Technologies Inc.\ (AIOps R\&D)}{Bengaluru, India}
{Software Development Engineer I \& II (ML/AI)}{Feb 2021 -- Jul 2025}
\resumeItemListStart
  \resumeItem{Led development of an agentic RAG diagnostics system from hackathon prototype to
  production deployment: LLM plans multi-step tool calls over network telemetry and verifies
  hypotheses via structured function calling -- autonomous root-cause analysis reducing mean
  resolution time ~70\%.}

  \resumeItem{Scaled OpenSearch telemetry ingestion pipelines from 50 to 2,000 events/sec (40x)
  via complete re-architecture (async I/O + Golang); deployed edge ML inference via ONNX Runtime
  (~40\% latency reduction); built SLA forecasting for 60+ production classifiers across 4 categories
  with automated evaluation and retraining.}

  \resumeItem{Patent (Filed): PCT/IN2022/058026 -- Unsupervised distributional thresholding for
  wireless connectivity anomaly detection; reduced manual troubleshooting ~75\%. Security-telemetry ML deployed in a production network-security environment (patented anomaly detection).}
\resumeItemListEnd
\resumeSubHeadingListEnd

%---------- PUBLICATIONS & PATENTS ----------
\section{Publications \& Patents}
\resumeSubHeadingListStart
\resumeProjectHeading
{\textbf{Patent (Filed):} PCT/IN2022/058026 -- Distributional thresholding for connectivity anomaly detection.}{}
\resumeProjectHeading
{\textbf{IEEE ICAIA 2026:} ML + Memory Forensics for IoT Cybersecurity (co-author).}{}
\resumeSubHeadingListEnd

%---------- TECHNICAL SKILLS ----------
\section{Technical Skills}
\footnotesize
\textbf{AI / ML:} PyTorch, Hugging Face Transformers, Deep Learning, Transformers, NLP, Computer Vision,
  Model Training \& Evaluation, Scikit-Learn, NumPy, Pandas, Experimentation \& Ablation Design \\[1pt]
\textbf{Agentic AI:} LangChain, LangGraph, Multi-Agent Frameworks, Tool/Function Calling, Agentic RAG,
  Planning \& Reflection, Long-Horizon Task Decomposition, Structured Generation (Pydantic) \\[1pt]
\textbf{Inference \& Deployment:} vLLM, Inference Optimization (ONNX Runtime), Latency/Throughput Engineering,
  Model Serving, Distributed Training (PyTorch DDP, Mixed Precision), Docker, Linux, Azure \\[1pt]
\textbf{Data Pipelines:} Python, Go, SQL/PostgreSQL, OpenSearch/Elasticsearch, Redis, Data Curation \&
  Quality Control, ETL, FAISS \\[1pt]
\textbf{Security ML:} Adversarial Robustness (PGD, Poisoning, Membership Inference -- coursework),
  Anomaly Detection (patented), Memory Forensics (published) \\[1pt]
\textbf{Reinforcement Learning:} PPO, DQN, Reward Shaping, Curriculum Learning, Self-Play (academic projects)

%---------- EDUCATION ----------
\section{Education}
\resumeSubHeadingListStart
\resumeSubheading
{Georgia Institute of Technology}{Atlanta, GA}
{\textbf{M.S.\ in Computer Science} (Machine Learning Specialization), \textbf{GPA: 4.0/4.0}}{Aug 2024 -- Dec 2026}
\vspace{2pt}
\begin{itemize}[leftmargin=0.20in, topsep=2pt, itemsep=0pt]
  \item\footnotesize{\textbf{Coursework:} Deep Learning, Deep RL, ML Security, Large \& Vision Language Models, Agentic AI, Systems for AI}
\end{itemize}
\vspace{-2pt}
\resumeEducationCompact
{Maharaja Surajmal Institute of Technology}{New Delhi, India}
{B.Tech.\ in Information Technology, CGPA: 8.8/10.0}{Aug 2016 -- Aug 2020}
\resumeSubHeadingListEnd

\end{document}
